\section{The same questions, two answers}
\label{sec:sameq}

The whole book revolves around five questions. The first part answered them in software, the
second part let the hardware do it:

\begin{booktable}{@{}lLL@{}}
\thead{Question} & \thead{Your own protocol} & \thead{CAN} \\ \midrule
Where does the message begin? & SOF, two bytes & The SOF bit; seven recessive bits as an end
  marker \\
How long is it? & LEN, one byte & DLC, four bits, 0--8 data bytes \\
Who is it for? & DST and SRC & \textbf{No address.} The identifier describes the content \\
Did it arrive intact? & A summing check\-sum & CRC-15, plus bit stuffing and format rules \\
Did it arrive at all? & ACK, timeout, retransmission, duplicate detection & One ACK bit inside the
  frame \\
\end{booktable}

Being able to fill in that table from memory, and to justify every row, is roughly what the
written test requires of the first half.

\section{What the written test covers}
\label{sec:exam}

\subsection{Part A --- your own protocol}

Frame design and checksum calculation by hand. Parsing as a state machine. Addressing, routing and
sequence numbers. Byte streams to be interpreted and analysed. The error model: dropped bytes,
corrupt data, delay. The chain ACK $\rightarrow$ timeout $\rightarrow$ retransmission
$\rightarrow$ duplicates.

\subsection{Part B --- CAN and the driver}

CAN IDs, priority, DLC and arbitration. Which mechanisms the hardware performs, and where. The
register map: what every register means, and how every status bit is cleared. The data byte order
across \code{TX_DATA_LO}/\code{HI}. The SPI transaction: the command byte, the four data bytes,
the \code{SS} rules, and encoding and decoding a transaction by hand. The layered architecture.
The three limitations.

The register map is handed out with the test --- the indices do not have to be memorised.
\emph{Reading} the map and knowing what it means is all the more necessary.
